% Data and Methodology

\subsection{Census Data}

We use decennial census data from three redistricting cycles:

\begin{table}[h]
\centering
\begin{tabular}{lrrr}
\toprule
\textbf{Year} & \textbf{Census Tracts} & \textbf{Total Population} & \textbf{Districts} \\
\midrule
2000 & 66,304 & 281,421,906 & 435 \\
2010 & 74,002 & 308,745,538 & 435 \\
2020 & 85,331 & 331,449,281 & 435 \\
\bottomrule
\end{tabular}
\caption{Census data summary across three decades}
\end{table}

Census tracts provide the finest-grained units with stable boundaries suitable for temporal comparison. Population data comes from the decennial census PL 94-171 redistricting files.

\textbf{Apportionment}: Congressional seats are allocated among states using the Huntington-Hill method:
\begin{itemize}
    \item 2000: Baseline allocation (435 seats across 50 states)
    \item 2010: 8 states gained, 10 states lost seats
    \item 2020: 6 states gained, 7 states lost seats
\end{itemize}

\subsection{Algorithmic Redistricting}

We apply recursive bisection using METIS 5.1.0 with \textit{identical parameters} across all three years:

\begin{itemize}
    \item \textbf{Algorithm}: Recursive bisection (graph partitioning)
    \item \textbf{Population balance}: $\pm0.5\%$ of target
    \item \textbf{Edge weights}: $\alpha = 5.0$ for minority-minority edges (maintaining communities of interest)
    \item \textbf{Contiguity}: Enforced via graph connectivity
    \item \textbf{No partisan data}: Algorithm has no access to voter registration or election results
\end{itemize}

\textbf{Why consistency matters}: Using identical parameters isolates temporal effects (population change, geographic shifts) from methodological variations. Any differences in outcomes across years reflect genuine demographic trends, not algorithmic artifacts.

\subsubsection{Voting Rights Act Considerations}

\textbf{Important limitation}: This algorithmic approach generates districts based solely on population balance and geographic compactness, without considering Voting Rights Act (VRA) Section 2 requirements for majority-minority districts. Real-world implementation would require VRA-compliant post-processing to ensure:

\begin{enumerate}
    \item \textbf{Gingles preconditions}: Where racial/ethnic minorities are sufficiently large and geographically compact to constitute a majority in a district, and where voting is racially polarized, Section 2 may require creating majority-minority districts \citep{gingles1986}.
    \item \textbf{Retrogression}: Districts must not diminish minority voting strength relative to existing plans \citep{beer1976}.
    \item \textbf{Proportionality consideration}: While not requiring strict proportional representation, VRA analysis considers whether minority groups have opportunities to elect candidates of choice proportional to their population \citep{johnson1997}.
\end{enumerate}

\textbf{Legal tension}: Redistricting authorities face a difficult balance: the VRA requires considering race when necessary to prevent vote dilution \citep{gingles1986}, but the Equal Protection Clause prohibits using race as the "predominant factor" in districting \citep{shaw1993}. In \textit{Shaw v. Reno}, the Supreme Court held that districts drawn predominantly on racial grounds, subordinating traditional redistricting principles like compactness, may constitute unconstitutional racial gerrymandering. This creates a narrow corridor: mapmakers must consider race enough to comply with Section 2, but not so much as to violate Equal Protection. Algorithmic approaches cannot navigate this legal tension automatically—human judgment remains essential.

We frame our algorithmic districts as a \textit{pre-VRA baseline}—a starting point demonstrating what purely geometric optimization produces before legal mandates are applied. In practice, redistricting authorities would need to adjust these algorithmically-generated boundaries to comply with VRA requirements, potentially reducing geometric compactness in favor of minority representation. Our longitudinal analysis shows how this baseline evolves over time, providing context for evaluating both enacted districts and VRA-compliant alternatives.

\subsection{Metrics}

\subsubsection{Compactness}

We measure compactness using the Polsby-Popper score:
\[
PP = \frac{4\pi \cdot \text{Area}}{\text{Perimeter}^2}
\]
PP ranges from 0 (infinitely elongated) to 1 (perfect circle). We compute PP for all 435 districts in each year.

\subsubsection{Geographic Stability}

To measure boundary stability across census cycles, we compute Intersection-over-Union (IoU):
\[
\text{IoU}(D_t, D_{t+10}) = \frac{|D_t \cap D_{t+10}|}{|D_t \cup D_{t+10}|}
\]
where $D_t$ is a district in year $t$ and $D_{t+10}$ is the geographically closest district 10 years later. IoU ranges from 0 (no overlap) to 1 (identical boundaries).

\subsubsection{Enacted Comparison}

We compare algorithmic districts to enacted districts using shapefiles from Dave's Redistricting App (DRA), which provides standardized enacted district boundaries for all three census cycles.

\subsection{Statistical Analysis}

We perform paired t-tests to assess compactness changes within states across years, ANOVA to compare commission vs non-commission states, and correlation analysis to examine relationships between population growth and district characteristics.
